\documentclass[aspectratio=169,12pt]{beamer}
\usepackage{fontspec}
\usepackage[portuguese]{babel}
\usepackage{xcolor,listings,tikz,booktabs,amsmath,amssymb,graphicx,multicol,array,colortbl}
\usetikzlibrary{arrows.meta,positioning,shapes.geometric,calc}

\definecolor{navy}{HTML}{0D1B3E}
\definecolor{gold}{HTML}{F5A623}
\definecolor{qgreen}{HTML}{1E8B4C}
\definecolor{qblue}{HTML}{1A6EAE}
\definecolor{codebg}{HTML}{EEF3F8}

\usetheme{default}
\useinnertheme{circles}
\setbeamertemplate{navigation symbols}{}
\setbeamertemplate{footline}{%
  \hfill{\color{gray!60}\scriptsize\insertframenumber/\inserttotalframenumber}\hspace{.6em}\vspace{.4em}}

\setbeamercolor{frametitle}{bg=navy,fg=white}
\setbeamercolor{structure}{fg=navy}
\setbeamercolor{alerted text}{fg=gold}
\setbeamercolor{item}{fg=gold}
\setbeamercolor{subitem}{fg=navy}
\setbeamercolor{block title}{bg=qblue,fg=white}
\setbeamercolor{block body}{bg=qblue!8,fg=black}
\setbeamercolor{block title alerted}{bg=gold,fg=navy}
\setbeamercolor{block body alerted}{bg=gold!15,fg=navy}
\setbeamercolor{block title example}{bg=qgreen,fg=white}
\setbeamercolor{block body example}{bg=qgreen!10,fg=black}
\setbeamerfont{frametitle}{size=\large,series=\bfseries}
\setbeamerfont{block title}{series=\bfseries}

\setbeamertemplate{title page}{%
  \begin{tikzpicture}[remember picture,overlay]
    \fill[navy](current page.north west)rectangle(current page.south east);
    \fill[gold]([yshift=.7cm]current page.south west)rectangle(current page.south east);
    \node[anchor=south,yshift=.73cm,navy,font=\tiny\bfseries]at(current page.south)
      {INTENSIVÃO QUANT AI \textbullet{} IMPACT UFSCAR};
  \end{tikzpicture}
  \vspace{.65cm}
  \begin{center}
    {\color{gold}\scriptsize\bfseries INTENSIVÃO QUANT AI \textbullet{} IMPACT UFSCAR}\\[.4cm]
    {\color{white}\LARGE\bfseries\inserttitle}\\[.25cm]
    {\color{gold!85}\normalsize\insertsubtitle}\\[.5cm]
    {\color{white!70}\small\insertauthor}\\[.1cm]
    {\color{white!50}\footnotesize\insertdate}
  \end{center}}

\lstdefinestyle{python}{language=Python,
  basicstyle=\ttfamily\scriptsize\color{qblue},
  keywordstyle=\color{navy}\bfseries,
  stringstyle=\color{qgreen!80!black},
  commentstyle=\color{gray!70}\itshape,
  backgroundcolor=\color{codebg},
  frame=l,framerule=2pt,rulecolor=\color{gold},
  xleftmargin=1em,breaklines=true,showstringspaces=false,tabsize=4,
  extendedchars=false,inputencoding=ascii}
\lstset{style=python}

\author{Felipe Maldonado\\{\scriptsize VP Diretoria Quant~\textbullet{}~Impact UFSCAR}}
\date{Intensivão Quant AI 2025}
\title{Aula 03 --- EDA}
\subtitle{Análise Exploratória de Dados Financeiros}

\begin{document}

\begin{frame}[plain]
\titlepage
\end{frame}

\begin{frame}[fragile]{Nesta Aula}
\begin{columns}[T]
  \column{0.5\textwidth}
\begin{block}{Teoria (30 min)}
\begin{itemize}
  \item Distribuições de retorno: normalidade vs fat tails
  \item Skewness, kurtosis e o Teorema do Limite Central
  \item Correlações: estáticas e rolling
  \item Estacionaridade e o teste ADF
\end{itemize}
\end{block}
  \column{0.47\textwidth}
\begin{block}{Código (70 min)}
\begin{itemize}
  \item Histograma + QQ-plot para diagnóstico de normalidade
  \item Scatter risk-return por ativo
  \item Heatmap de correlação com clusterização
  \item Correlação rolling de 12 meses
  \item Teste ADF: quais ativos são estacionários?
\end{itemize}
\end{block}
\end{columns}
\end{frame}

\begin{frame}[fragile]{Distribuição de Retornos Financeiros}
\begin{columns}[T]
  \column{0.5\textwidth}

\textbf{Hipótese Normal (incorreta):}
\[ r_t \sim \mathcal{N}(\mu, \sigma^2) \]

\textbf{Realidade empírica:}
\begin{itemize}
  \item \textbf{Kurtosis} $> 3$: caudas mais pesadas que Normal
  \item \textbf{Skewness} $\neq 0$: assimetria, geralmente negativa
  \item \textbf{Autocorrelação} de $r_t$: próxima de zero
  \item \textbf{Autocorrelação} de $|r_t|$: positiva (volatility clustering)
\end{itemize}
\vspace{.3em}
\begin{alertblock}{Implicação para risco}
VaR e CVaR calculados com Normal \textbf{subestimam} eventos extremos.
Distribuições: $t$-Student, GED, misturas de normais.
\end{alertblock}
  \column{0.47\textwidth}
\begin{block}{Stylized Facts (Cont, 2001)}
\begin{itemize}
  \item Heavy tails: kurtosis $\approx$ 4--10 em ações
  \item Gain/Loss Asymmetry: quedas mais abruptas que altas
  \item Volatility Clustering: $|\varepsilon_t|$ correlacionado
  \item Leverage Effect: $\text{Corr}(r_t, \sigma_{t+k}^2) < 0$
  \item Long memory em volatilidade (GARCH)
\end{itemize}
\end{block}\vspace{.3em}
\begin{exampleblock}{QQ-Plot}
Quantis empíricos vs Normal: pontos nas caudas \textbf{acima} da linha $\Rightarrow$ fat tails.
\end{exampleblock}
\end{columns}
\end{frame}

\begin{frame}[fragile]{Correlações e Estacionaridade}
\begin{columns}[T]
  \column{0.5\textwidth}
\begin{block}{Matriz de Correlação}
\begin{itemize}
  \item Correlações cross-sectional: como os ativos se movem juntos
  \item Hierarquical clustering: agrupa setores similares
  \item Em crises: correlações convergem para 1 (diversificação falha)
  \item Rolling 12m: evolução temporal das relações
\end{itemize}
\end{block}\vspace{.3em}
\begin{alertblock}{Correlação Espúria}
Séries não-estacionárias têm alta correlação por tendência compartilhada — não por relação causal.
\end{alertblock}
  \column{0.47\textwidth}
\begin{block}{Teste ADF}
\begin{itemize}
  \item $H_0$: série tem raiz unitária (não-estacionária)
  \item $H_1$: série é estacionária
  \item Retornos $r_t$: geralmente estacionários ($p < 0{,}05$)
  \item Preços $P_t$: geralmente não-estacionários (random walk)
\end{itemize}
\end{block}\vspace{.3em}
\begin{exampleblock}{Por que importa?}
Modelos de regressão em séries não-estacionárias produzem regressão espúria. Sempre transforme preços em retornos antes de modelar.
\end{exampleblock}
\end{columns}
\end{frame}

\begin{frame}[fragile]{O que Vamos Construir}
\begin{columns}[T]
  \column{0.52\textwidth}
\begin{block}{Funções a implementar}
\begin{itemize}
  \item \texttt{plot\_histograma\_qqplot(ret)} $\to$ diagnóstico de normalidade
  \item \texttt{plot\_risk\_return(ret)} $\to$ scatter por ativo
  \item \texttt{plot\_heatmap\_correlacao(ret)} $\to$ matriz com cluster
  \item \texttt{plot\_rolling\_corr(ret, janela=12)} $\to$ evolução temporal
  \item \texttt{testar\_estacionaridade(ret)} $\to$ ADF por ativo
\end{itemize}
\end{block}
  \column{0.45\textwidth}
\begin{block}{Entradas (parquets)}
\begin{itemize}
  \item \texttt{retornos\_mensais\_limpo.parquet}
  \item \texttt{retornos\_diarios\_limpo.parquet}
\end{itemize}
\end{block}
\vspace{.4em}
\begin{exampleblock}{Saídas (parquets)}
\begin{itemize}
  \item Gráficos PNG para o relatório
  \item DataFrame com p-valores do ADF
\end{itemize}
\end{exampleblock}
\end{columns}
\end{frame}

\begin{frame}[fragile]{Referências}
\begin{multicols}{2}
\tiny
\begin{itemize}
  \item \textbf{Cont, R.} (2001). Empirical properties of asset returns. \textit{Quantitative Finance}, 1(2), 223--236.
  \item \textbf{Mandelbrot, B.} (1963). The variation of certain speculative prices. \textit{Journal of Business}, 36(4), 394--419.
  \item \textbf{Engle, R.} (1982). Autoregressive conditional heteroscedasticity. \textit{Econometrica}, 50(4), 987--1007.
  \item \textbf{Dickey, D. \& Fuller, W.} (1979). Distribution of the estimators for autoregressive time series. \textit{JASA}, 74(366), 427--431.
\end{itemize}
\end{multicols}
\end{frame}


\end{document}
